% TOP_PROBLEM: {"id": 3232, "title": "Partial List Coloring", "category_id": 3, "category": {"id": 3, "name": "graph_theory", "display_name": "Graph Theory", "description": "Problems involving graphs, networks, and their properties.", "slug": "graph-theory", "order_index": 3, "created_at": "2026-07-31T15:26:25.671Z"}, "status": "open", "problem_number": "OPG-788", "set_id": 12}
% FIRST_POSTED: 2026-10-01
% ATTEMPT_STATUS: unresolved
% MODEL: GPT 6 Astra Ultra
% UnsolvedMath ID 3232; source catalogue code OPG-788
% !TeX program = lualatex
% Research attempt dated 2026-10-01; canonical statement preserved.
\documentclass[11pt,a4paper]{article}
\usepackage[margin=25mm]{geometry}
\usepackage{fontspec}
\setmainfont{lmroman10-regular.otf}[BoldFont=lmroman10-bold.otf,
 ItalicFont=lmroman10-italic.otf,BoldItalicFont=lmroman10-bolditalic.otf]
\usepackage{amsmath,amssymb,mathtools}
\usepackage{unicode-math}
\setmathfont{latinmodern-math.otf}
\AtBeginDocument{\let\setminus\smallsetminus}
\usepackage{xurl,hyperref}
\hypersetup{colorlinks=true,urlcolor=blue}
\setcounter{secnumdepth}{0}
\setlength{\parindent}{0pt}
\setlength{\parskip}{5pt}
\setlength{\emergencystretch}{3em}
\begin{document}
\section*{Partial List Coloring}\label{3232}
{\small UnsolvedMath ID 3232 · OPG-788 · Graph Theory · OpenGarden}
\subsection{Definitions and mathematical statement}
Let $G=(V,E)$ be a finite nonempty simple undirected graph. A list assignment
$L$ associates a finite set $L(v)$ of available colours with each vertex $v$.
An $L$-colouring of an induced subgraph $G[U]$, where $U\subseteq V$, is a
map $c:U\to\bigcup_{v\in V}L(v)$ such that $c(v)\in L(v)$ and
$c(u)\ne c(v)$ whenever $uv\in E$ has both ends in $U$. Put
\[
 \lambda_L(G)=\max\{|U|:G[U]\text{ has an }L\text{-colouring}\},
 \qquad
 \lambda_t(G)=\min_{|L(v)|=t\ (v\in V)}\lambda_L(G).
\]
The minimum ranges over arbitrary colour sets and arbitrary overlaps between
lists. Although there are infinitely many names for colours, the possible
values are integers between $0$ and $|V|$, so this minimum is well defined.

The list chromatic number $\chi_\ell(G)$ is the least positive integer $s$
for which every assignment of $s$ colours per vertex colours all of $G$.
The conjecture asks whether, for every such graph and every pair of integers
$1\le r\le s\le\chi_\ell(G)$,
\[
                 \frac{\lambda_r(G)}r\ge\frac{\lambda_s(G)}s.
\]
The quantity compares the guaranteed number of colourable vertices per
available colour. The optimizing subset may depend on the entire list
assignment; no particular subset must work for every assignment. At the
endpoint $s=\chi_\ell(G)$ the assertion reads
$\lambda_r(G)\ge r|V|/\chi_\ell(G)$, since all vertices can then be coloured.
\subsection{Short English statement}
Does the worst-case number of vertices that can be list-coloured, divided by the list size, decrease as the lists become larger?
\subsection{Research attempt}
\paragraph{Scope and method.}
This research attempt by GPT-6 Astra Ultra is dated 1 October 2026.
It preserves the catalogue statement and distinguishes elementary proofs
below from cited prior work. No novelty claim or formal proof-assistant
verification is made. The final paragraph states the precise remaining scope.
\paragraph{Status and a precise implication.}
The historical monotonicity claim is stronger than the endpoint assertion
in UnsolvedMath 3231. Noel [R1] reports a graph with $n=14$,
$\chi_\ell=3$ and $\lambda_2\le9$. It follows that
\[
           \frac{\lambda_2}{2}\le\frac92
           <\frac{14}{3}=\frac{\lambda_3}{3}.
\]
Thus that primary counterexample also refutes this catalogue formulation.
We retain its definition and investigate which comparison principles do
remain valid. None of the following positive restricted results rescues
the universal assertion at the displayed parameters.

\paragraph{Monotonicity without normalization.}
If $r\le s$, then $\lambda_r\le\lambda_s$. Indeed, for every $s$-list
assignment restrict each list to any $r$ elements. The restricted lists
permit a partial colouring of at least $\lambda_r$ vertices, which is
also a valid colouring from the original lists. Taking the minimum over
the original $s$-assignments proves the claim. This controls the numerator
only; dividing by the increasing list size can reverse comparisons.

\paragraph{Disjoint palettes prove subadditivity.}
Choose assignments attaining $\lambda_a$ and $\lambda_b$, and rename
their colours so that their palettes are disjoint. Their vertexwise
unions are $(a+b)$-lists. Every partial colouring from these unions
splits into two vertex sets according to which palette its colour uses.
Each set is a proper partial colouring from its corresponding original
assignment, so the total has size at most $\lambda_a+\lambda_b$.
Since $\lambda_{a+b}$ minimizes over all $(a+b)$-assignments, we obtain
\[
                         \lambda_{a+b}\le\lambda_a+\lambda_b.
\]
The minima exist because their possible values form a finite nonempty
set of integers, even though the colour names range over arbitrary sets.
Interactions between the two vertex sets cannot invalidate this upper
bound; if anything, such interactions further restrict a colouring.

\paragraph{Divisibility and the singleton endpoint.}
Taking $k$ disjoint renamed copies of an assignment attaining $\lambda_r$
gives $\lambda_{kr}\le k\lambda_r$. Hence the desired normalized comparison
does hold whenever $s=kr$ and $s\le\chi_\ell(G)$. In particular, when
$r$ divides $q=\chi_\ell(G)$ it implies $\lambda_r\ge rn/q$ because
$\lambda_q=n$. This is a genuine endpoint theorem on a restricted set
of parameters, and it excludes the nondivisible pair $r=2,q=3$ above.

Also $\lambda_1=\alpha(G)$. Any maximum independent set can always use
its singleton lists, and assigning one common colour everywhere prevents
more than $\alpha(G)$ vertices from being coloured. Applying subadditivity
repeatedly yields $\lambda_s\le s\alpha(G)$, proving the normalized
comparison for $r=1$ against every $s$.

\paragraph{A full positive graph class.}
For a disjoint union of cliques with orders $m_i$,
$\lambda_t=\sum_i\min(t,m_i)$. For the lower bound choose at most $t$
vertices in each clique and apply Hall's condition to their $t$-element
lists. Identical lists on every vertex attain the upper bound. Dividing
by $t$ gives $\lambda_t/t=\sum_i\min(1,m_i/t)$, a sum of nonincreasing
functions of positive $t$. Thus the full comparison holds for this class,
including unequal clique sizes, rather than only for divisible indices.

\paragraph{Verification and the logical limit of subadditivity.}
The root batch script checks the clique formula and normalized comparisons
for all sequences of up to four clique sizes between one and six.
It also confirms the finite bad-list count used in [R1]. Finite palette
search is not used to prove a universal choice number.
Pure subadditivity cannot prove the missing ratios: the abstract integer
sequence $0,3,4,7$ is nondecreasing and subadditive on indices through
three, but $4/2<7/3$. This sequence is only a logical diagnostic, not a
claim that a graph realizes those values.

\paragraph{Final status.}
\textbf{UNRESOLVED as a new-proof attempt.} The cited counterexample
transfers to the historical universal assertion by the exact displayed
inequality. Divisibility, the singleton endpoint, and all clique unions
are proved positive cases. No independent new disproof or optimal
replacement monotonicity theorem is claimed.

\subsection{Sources}
\begin{itemize}
\item[S1] ULAM.AI and the UnsolvedMath Contributors, supplied problems.json, record 3232 (OPG-788), OpenGarden collection. Catalogue identity and source transcription; CC BY 4.0.
\url{https://huggingface.co/datasets/ulamai/UnsolvedMath}.
\item[S2] Open Problem Garden, Partial List Coloring. Original problem and discussion; the mathematical formulation below is expanded from the supplied source record.
\url{http://www.openproblemgarden.org/op/partial_list_coloring_0}.
\item[R1] Jonathan A. Noel, \emph{The Partial List Colouring
Conjecture is False}, arXiv:2609.23291v1, 20 September 2026.
Primary full text checked 1 October 2026.
\url{https://arxiv.org/html/2609.23291v1}.
\item[R2] Moharram Iradmusa, \emph{A Note on Partial List Colorings},
arXiv:0805.3277. Primary abstract checked 1 October 2026; historical
source for related partial-list comparisons.
\url{https://arxiv.org/abs/0805.3277}.
\end{itemize}
\end{document}
